% RECOVERY32: reconstructed from actual_sphere_cabling_correction31.md.
\statement{A.98}{Proved}{The original cabling arrows determine the full correction}
Let \(\alpha=(-1,\varphi)\in\widehat{GT}\) be the actual involution,
with canonical \(u\in\widehat F'\) satisfying \(u=\varphi\alpha(u)\).
Use A.96--A.97, retaining \(f=A_{34}\) for the braid generator.
There are actual words
\[
 p=\varphi(bc,f),\qquad z_u=u(bc,f)
       \quad\hbox{in }(\widehat N)',
 \tag{A.98.1}
\]
such that \(z_u=p\alpha(z_u)\) and
\[
 \gamma=\operatorname{Inn}(p)\alpha
 =\operatorname{Inn}(z_u)\alpha\operatorname{Inn}(z_u^{-1}).
 \tag{A.98.2}
\]
The free factor \(\widehat F(A,S)\le\widehat N\), \(S=CP\),
is a literal copy of the original canonical system:
\[
 \gamma(A)=A^{-1},\qquad
 \gamma(S)=\varphi(A,S)^{-1}S^{-1}\varphi(A,S).
 \tag{A.98.3}
\]
If t is the unique transporter in this free factor with
\(\gamma(L)=t^{-1}L^{-1}t\), \(L=SA\), and zero L-coordinate
in the basis \((L,A)\), then the actual A.97 correction is
\[
 h_0=tp,\qquad t_{\rm ab}=-A,
       \qquad (h_0)_{\rm ab}=-A.
 \tag{A.98.4}
\]
In particular h0 is not an alpha-coboundary in \(\widehat N\),
and the corresponding section distinction persists in every actual
relative pro-two quotient of A.97.

\begin{proof}
Use the actual objects and underlying-identity associativity arrows
of A.36:
\[
 b_0=(((12)3)4),\quad b_4=((12)(34)),\quad b_3=(1(2(34))),
 \quad D_1:b_0\to b_4,\quad D_2:b_4\to b_3.
\]
Their original source-based convention is
\[
 \alpha(D_1)=D_1\varphi(bc,f)=D_1p,
 \qquad \alpha(D_2)=D_2\varphi(a,ce).
 \tag{A.98.5}
\]
At b0, the block twist bc is the quotient of the commuting,
literally inverted twists abc and a. At b4, f is an internally
grouped pair twist and is literally inverted. Transport through D1
therefore gives
\(\alpha(bc)=(bc)^{-1}\), \(\alpha(f)=p^{-1}f^{-1}p\).
This is exactly the original normalized two-generator action.

In the sphere presentation \(bc=f^{-1}R^{-1}\), \(R=PA\).
Because the base action is IA,
\[
 \widehat G/(\widehat N)'
       =(\widehat N)_{\rm ab}\times\widehat B,
 \qquad (\widehat N)_{\rm ab}=\widehat{\mathbb Z}^{3}.
\]
The arguments bc and f commute in this quotient. Both parameters
are commutator words, proving (A.98.1). Evaluating the actual norm
equation gives \(z_u=p\alpha(z_u)\), and hence (A.98.2), with
the displayed order. This is a proved conjugacy, unlike an arbitrary
inner twist of the original involution.

The words \((P,A,S)\) form a free basis of N, and f fixes A,S
literally by A.96. Thus the group on A,S,f is their direct product
with the f-factor. The original arguments are
\(a=Af\), \(ce=Sf^{-1}\). Every profinite commutator word
evaluated at these arguments equals its evaluation at A,S: the
central f-coordinate is zero, first for words and then by continuity.
At b4 the action is gamma, which inverts a and f literally.
The second arrow in (A.98.5) gives the second formula in
(A.98.3), and the first follows from \(A=af^{-1}\).

The original third-cusp condition in this literal rank-two copy
gives t. The primitive centralizer of L makes its zero-L-coordinate
normalization unique. From (A.98.2), \(h_0p^{-1}\) is a
transporter for the same equation. Since p belongs to
\((\widehat N)'\), its L-coordinate remains zero. The free-factor
normalization agrees with the normalization in \((L,A,C)\);
uniqueness in \(\widehat N\) proves \(h_0=tp\).

For the integral abelianization calculation, use the class-two
Heisenberg quotients of \(F(A,S)\) over every \(\mathbb Z/n\).
The commutator word \(\varphi(A,S)\) is central there, so
\(\gamma(L)=S^{-1}A^{-1}\). The element \(A^{-1}\) is a
transporter, since \(AL^{-1}A^{-1}=S^{-1}A^{-1}\).
Consequently tA centralizes L in every such quotient. Zero
L-coordinate means \(t_{\rm ab}=(r,0)\) in the (A,S)-basis.
The commutator of tA with L has coefficient \(r+1\); vanishing
modulo every n gives \(r=-1\) in \(\widehat{\mathbb Z}\).
This proves all of (A.98.4), without a division or finite census.

Alpha acts as minus identity on the compact abelianization of N.
Every coboundary \(q\alpha(q)^{-1}\) consequently has even
abelian coordinates, while h0 has odd A-coordinate. The sections
alpha and \(h_0\alpha\) are therefore not N-conjugate. In every
relative quotient of A.97, the embedded \(N^{(2)}\) retains that
A-character modulo two. It extends by zero on the base, because
the action is IA, and is alpha-invariant. Thus it vanishes on
every coboundary and is one on h0, proving the same distinction
in the actual relative group.

Finally the B-action on \(N^{(2)}\) factors through \(B^{(2)}\):
on every characteristic finite two-group quotient its image lies
in the Frattini automorphism kernel, a two-group by A.96.
Taking all those quotients proves the factorization. This does not
erase the odd base head, identify the negative transporter with
a regular one, or decide the nonlinear cusp equation. It records
the exact same correction and the actual limit of this kernel action.
\end{proof}
